% New file while the preceding checkpoint is under visual review.
\paragraph{C.1i (Proved: all non-top cases of $A'$ and $A''$ fail universally).}
Fix integers $c<n<2c$. There is a specified smooth projective
$n$-fold $X$ over the algebraic closure of a rational function
field over $\mathbb Q$, with a very ample class $L$, and a class
\[
 0\ne z\in G^c(X),\qquad
 z\notin L A^{c-1}_{\mathrm{alg}}(X).
\]
After adjoining coefficients for a geometric generic hyperplane
$Y\subset X$, the image of $z$ is outside
$i_*A^{c-1}_{\mathrm{alg}}(Y)$, and it is nonzero in
$A^c_{\mathrm{alg}}(X\setminus Y)$. These statements remain true
after every embedding of the algebraically closed defining field
into $\mathbb C$.

In particular, in the letter's range $2c>n$, with $1\le c\le n$,
the assertions $A'$ and $A''$ hold for every smooth projective
complex $n$-fold and smooth hyperplane section exactly when $c=n$.
Here the index in $A'$ is $i=c-1$; its range $2i\ge n-1$ is
exactly $2c>n$. Groups with $c>n$ vanish for dimension reasons.

\emph{Construction and proof.}
We use the split quadric
\[
 Q=\{x_0y_0+x_1y_1+\cdots+x_cy_c=0\}
            \subset\mathbb P^{2c+1}_{\mathbb Q},
       \qquad h=c_1(\mathcal O_Q(1)),
\]
and its two maximal linear spaces
\[
 \Pi_+=\{x_0=x_1=\cdots=x_c=0\},\qquad
 \Pi_-=\{y_0=x_1=\cdots=x_c=0\}.
\]
Put $Z=[\Pi_+]-[\Pi_-]\in\mathrm{CH}^c(Q)$.
The linear section $x_1=\cdots=x_c=0$ is the reduced union
$\Pi_+\cup\Pi_-$, so $h^c=[\Pi_+]+[\Pi_-]$.
Both spaces have degree one. The normal-bundle sequence
\[
 0\longrightarrow N_{\Pi_+/Q}\longrightarrow
       \mathcal O_{\mathbb P^c}(1)^{\oplus(c+1)}
       \longrightarrow\mathcal O_{\mathbb P^c}(2)
       \longrightarrow0
\]
is exact, since the displayed quadric is smooth. It gives
\[
 [\Pi_+]^2=[t^c]\frac{(1+t)^{c+1}}{1+2t}
       =\frac{1+(-1)^c}{2},\qquad
 \deg Z^2=2(-1)^c\ne0.
\]
Thus $\operatorname{cl}(Z)\ne0$. Weak Lefschetz and duality give
$H^{2c-2}(Q,\mathbb Q)=\mathbb Qh^{c-1}$ and
$\dim H^{2c+2}(Q,\mathbb Q)=1$. Equal degrees imply
$h\operatorname{cl}(Z)=0$. Consequently
\begin{equation}\label{eq:quadric-primitive-obstruction}
 \operatorname{cl}(mZ-ha)\ne0
 \quad\text{for every }m\ne0\text{ and }
                 a\in\mathrm{CH}^{c-1}(Q)_{\mathbb Q}.
\end{equation}
Indeed, the class of $ha$ is a multiple of $h^c$ and is
orthogonal to $Z$, whereas $Z^2\ne0$.

Set $s=2c-n$ and $D=4c+2$. Choose a basis of
$V=H^0(Q,\mathcal O_Q(D))$ over $\mathbb Q$, let
$S=\mathbb P(V^{\oplus s})$, and let $X$ be the geometric
generic common zero locus of its $s$ sections. This is smooth
of dimension $n$ by Bertini. To specify the field and equations
without any unnamed generality choice, use the degree-$D$
monomials not divisible by $x_0y_0$ as the basis of
$\mathbb Q[x_0,y_0,\ldots,x_c,y_c]/(\sum x_jy_j)$.
In the first section set the coefficient of $x_0^D$ equal to one;
take every other coefficient of the $s$ sections to be an
independent indeterminate. Thus the field is
$K=\mathbb Q(t_1,\ldots,t_{s\dim V-1})$ and we work over
$\Omega=\overline K$. Write $j:X\hookrightarrow Q_\Omega$,
$L=j^*h$ and $z=j^*Z$.

We make the only deep input explicit. We use the effective
connectivity theorem of \citet[Theorem 3.13, p.~23 of the opened
author version]{cyNagel2002}, with the required condition that
its connectivity parameter be at most $\dim X$, as made explicit
in \citet[Theorem 5.11]{cyNagel2004}. The regularity constant
$m_Q=1$ for a smooth quadric is recorded in
\citet[p.~3 of the opened author version, immediately before
Theorem 2]{cyNagel1998}. Put its parameter equal to $n$;
then $\mu=n$ and $r=s-1$. Since all degrees equal $D$, condition
$(C)$ has left side at least $D>2c$, and each condition $(C_i)$
has left side
\[
  \sum_{\nu=i}^{r}D+iD=sD\ge D>1+2c+n-i.
\]
Thus all the hypotheses hold. For every smooth morphism
$T\to S^\circ$, where $S^\circ$ is the smooth-fibre open,
the theorem gives
\begin{equation}\label{eq:nori-effective-connectivity}
 H^k(Q\times T,X_T;\mathbb Q)=0\qquad(k\le2n).
\end{equation}
The parameter $S=\mathbb P H^0(Q,\mathcal O_Q(D)^{\oplus s})$
is precisely the one in that theorem. No ineffective degree
threshold or unspecified exceptional field is used here.

For clarity, we reconstruct the consequence of connectivity that
we need. It is the mechanism of
\citet[Theorem 1, pp.~349--350 and Sections 5--6]{cyNori1993}.
If $W\in\mathrm{CH}^c(Q)_{\mathbb Q}$ becomes algebraically
trivial on the geometric generic $X$, clear denominators and
spread the finite equivalence witness after a finite extension
of $\mathbb Q(S)$. After shrinking, it gives a smooth base
$T$ finite \emph{\'etale} over an open of $S^\circ$, a smooth
projective family of curves $C\to T$, a divisor $\delta$ of
degree zero on every connected fibre component, and a
correspondence $\Gamma\in\mathrm{CH}^c(X_T\times_T C)$ with
\[
       j_T^*W=(p_X)_*(\Gamma\cdot p_C^*\delta)
                   \quad\text{in }\mathrm{CH}^c(X_T)_{\mathbb Q}.
\]
A disjoint union of curve families combines finitely many
relations. Rational-equivalence witnesses also spread after
shrinking, so the displayed equality is an actual Chow equality
on the total space, not merely an equality at its fibres.
Base change the finite-type data to $\mathbb C$. Both
$T\to S^\circ$ and $C\to S^\circ$ are smooth. Because $c<n$,
\eqref{eq:nori-effective-connectivity} makes
\[
 H^{2c}(Q\times C,\mathbb Q)
         \longrightarrow H^{2c}(X_T\times_T C,\mathbb Q)
\]
surjective, and makes
$H^{2c}(Q\times T,\mathbb Q)\to H^{2c}(X_T,\mathbb Q)$
injective. Lift $\operatorname{cl}(\Gamma)$ to a class $\beta$
on $Q\times C$, and push $\beta\,p_C^*\operatorname{cl}(\delta)$
along the proper curve projection $Q\times C\to Q\times T$.
The resulting class restricts to $j_T^*\operatorname{cl}(W)$,
so injectivity identifies it with the pullback of
$\operatorname{cl}(W)$. Restriction to one fibre over $T$ is
zero, because a degree-zero divisor on a projective curve has
zero cohomology class. We conclude
\begin{equation}\label{eq:nori-geometric-generic-injectivity}
 j^*W=0\text{ modulo algebraic equivalence over }\Omega
       \quad\Longrightarrow\quad\operatorname{cl}(W)=0.
\end{equation}
This proof deals directly with a putative finite witness on the
specified geometric generic fibre; it does not assume that
analytic exceptional loci are defined over $\overline{\mathbb Q}$.

The class $z$ is homologically trivial. Indeed, for every
$b\in H^{2n-2c}(X,\mathbb Q)$, weak Lefschetz gives
$b=j^*b_0$ for a class on $Q$, since $2n-2c<n$.
The projection formula gives
\[
 \int_X j^*\operatorname{cl}(Z)\,b
       =D^s\int_Q\operatorname{cl}(Z)\,h^s b_0=0,
\]
because $s\ge1$ and $h\operatorname{cl}(Z)=0$.
Poincar\'e duality proves the assertion. Equation
\eqref{eq:nori-geometric-generic-injectivity} proves $z\ne0$
modulo algebraic equivalence.

We now prove the stronger obstruction $z\notin L A^{c-1}_{\mathrm{alg}}(X)$.
Suppose otherwise and choose a finite extension $K'/K$ over
which a cycle $\alpha$ and the finite algebraic-equivalence
witness for $z=L\alpha$ are both defined. Trace to $K$ and
write $e=[K':K]$. The trace relation is
\[
       e z=L\alpha_{\mathrm{tr}}
              \quad\text{modulo algebraic equivalence over }\Omega.
\]
The universal incidence $\mathcal X\subset Q\times S$ is a
projective bundle over $Q$: it is the projectivization of the
kernel of the surjective evaluation map
\[
 V^{\oplus s}\otimes\mathcal O_Q
       \longrightarrow\mathcal O_Q(D)^{\oplus s}.
\]
Its relative hyperplane class is pulled back from $S$ and
restricts to zero on its generic fibre. Close
$\alpha_{\mathrm{tr}}$ in $\mathcal X$ and use the
projective-bundle formula. Restriction to the generic fibre
therefore writes
$\alpha_{\mathrm{tr}}=j^*a$ for
$a\in\mathrm{CH}^{c-1}(Q)_{\mathbb Q}$.
We obtain $j^*(eZ-ha)=0$ modulo algebraic equivalence.
Equation \eqref{eq:nori-geometric-generic-injectivity} contradicts
\eqref{eq:quadric-primitive-obstruction}. This proves the claim.

Adjoin a geometric generic hyperplane to $X$ and apply C.1h.
The obstruction survives this algebraically closed extension:
any putative preimage and relation over the larger field spread
over a finite-type model and specialize to a point over $\Omega$,
contradicting the obstruction just proved. The same finite-data
argument proves persistence under embeddings into $\mathbb C$.
Finally, when $c=n$, degree identifies
$A^n_{\mathrm{alg}}(X)$ with $\mathbb Q$; a point of $Y$
pushes to a class of degree one. Hence $A'$ is onto and
localization gives $A^n_{\mathrm{alg}}(X\setminus Y)=0$.
This proves all the asserted universal cases.\hfill$\square$
